% Abstract: tradução do resumo (editaveis/resumo.tex), com os mesmos números.
\begin{resumo}[Abstract]
 \begin{otherlanguage*}{english}
   This work develops a prototype open source intelligence (OSINT) search
   tool for the Brazilian judicial police to investigate fraud in public
   procurement. The requirements came from an elicitation with an expert and
   were checked against the legislation and a systematic literature review
   (SLR 1). A second review (SLR 2) guided the development process. Both
   reviews are reported following PRISMA 2020 and SEGRESS. SLR 1 covers OSINT
   tools for law enforcement and intelligence agencies. Of the 1,928 records
   retrieved, 1,139 remained after duplicate removal, and 57 primary studies
   were included, 23 from the database search and 34 from snowballing. No
   study is from Brazil or any other Latin American country, none addresses
   fraud or corruption in public procurement, and only 7 explicitly address
   the chain of custody. SLR 2 covers the effects of
   assistants and agents based on large language models on software
   development. Of the 8,181 records, 4,323 remained after duplicate removal,
   and 130 reports, corresponding to 129 studies, were included. On
   productivity, 51 of the 89 studies that measured the effect report a
   positive direction, and the direction is more favorable when the effect is
   only perceived. On security, 14 of the 23 studies that measured the effect
   report a negative direction. A comparison with seven existing tools, based
   on their official documentation, found none that combines Brazilian open
   procurement data sources, fraud typologies and evidence preservation for
   the police inquiry. The requirements were written as 39 user stories,
   prioritized with the MoSCoW method. The open source prototype has passive
   collection only, a chain of custody with a SHA-256 hash and a timestamp for
   each capture, an append-only audit trail, ten fraud typologies based on
   explicit rules, whose alerts are leads that the investigator confirms or
   dismisses, and optional, local AI-based entity extraction. It was
   developed with AI agents, applying some of the candidate practices from
   SLR 2, and has 342
   automated back-end tests and 234 front-end tests. The review results are
   provisional, because steps performed by AI as a third opinion still await
   confirmation by the two reviewers. The prototype has automated tests with synthetic
   data and its connectors were checked against the real sources, but it has not yet been evaluated with police officers. The conclusion is that the literature and the tools
   analyzed do not cover the combination of the public procurement domain and
   evidence preservation, and that the prototype shows that this combination
   can be implemented.

   \vspace{\onelineskip}

   \noindent
   \textbf{Keywords}: open source intelligence. criminal investigation.
   public procurement. chain of custody. systematic review.
 \end{otherlanguage*}
\end{resumo}
